%-------------------------------------------------------------------------------
% This is a master's thesis for KIDS students.

% The template used in this article was provided by  Sofia Feist, 2020,
% and has been modified on this basis.
%------------------------------------------------------------------------------- 

% ------------------ DOCUMENT SETUP ------------------ 
% The document class defines the document type (report) and sets the font size (12pt)
\documentclass[12pt]{report}
\author{SRN:XXX}

% Inputs the Document Packages
\input{_Packages}

% Font
\setmainfont{Arial}

% Line Spacing
\setstretch{2.0}

% Controls how many subchapters the document can take
%  and how many of those will get put into the contents pages.
\setcounter{secnumdepth}{2}
\setcounter{tocdepth}{2}

% The folder path where the images will be uploaded
\graphicspath{{./Images/}} 

% Numeration Type for chapters and chapters (Roman I, II, II / Arabic 1, 2, 3)
% \renewcommand\thechapter{\Roman{chapter}}
\renewcommand\thechapter{\arabic{chapter}}

% Line Break Properties
\tolerance=1
\emergencystretch=\maxdimen
\hyphenpenalty=10000
\hbadness=10000

% Formatting Table of Contents/Lists titles
\renewcommand{\contentsname}{\normalfont\bfseries\LARGE{CONTENTS}}
\renewcommand{\listfigurename}{\normalfont\bfseries\LARGE{LIST OF FIGURES}}
\renewcommand{\listtablename}{\normalfont\bfseries\LARGE{LIST OF TABLES}}

% Title Formatting customization
\titleformat{\chapter}{\normalfont\bfseries\LARGE}{\thechapter.}{1em}{\MakeUppercase}
\titlespacing*{\chapter} {0pt} {0pt} {15pt} % left, before, after

\titleformat{\section}{\normalfont\bfseries\Large}{\thesection.}{1em}{}
% \titlespacing*{\section} {0pt} {0pt} {15pt}

\titleformat{\subsection}{\normalfont\bfseries\large}{\thesubsection.}{1em}{}
% \titlespacing*{\subsection} {0pt} {10pt} {10pt}

\titleformat{\subsubsection}{\normalfont\bfseries\normalsize}{}{1em}{}
% \titlespacing*{\subsubsection} {0pt} {10pt} {10pt}

% Set headheight and adjust topmargin
\setlength{\headheight}{14.5pt}
\addtolength{\topmargin}{-2.5pt}

% HEADER AND FOOTER
\pagestyle{fancy}  % Set Page Style (Header and Footer Style)
\fancyhf{}  % Clears the header and footer (from the default info)

% Header
\fancyhead[L]{your-moodle-code}
\fancyhead[C]{\thepage}
\fancyhead[R]{SRN:XXX}
% \renewcommand{\headrulewidth}{0pt}  % Removes the default Horizontal Line in Header


% Change figure numbering per chapter
\numberwithin{figure}{chapter}
\numberwithin{table}{chapter}
\numberwithin{equation}{chapter}

%  -------------------------------------------------
%  --------- The document starts from here --------- 
%  -------------------------------------------------
\begin{document}

% ------------------  TITLE PAGE -------------------
\begin{titlepage}
\begin{center}
    % UCL IMAGE
    \vspace*{-3.5cm}
    \makebox[\textwidth]{\includegraphics[width=\paperwidth]{Image/ucl_letterhead.png}}
    
    \vfill % Elastic empty space filler

    \vspace{1cm}
    % Title
    {\Huge\textbf{Dissertation Title\\
    Dissertation Subtitle\\}}

    \vspace{2cm}

    % Author
    {\Large\textbf{SRN: XXX\\}}
    % Date
    {\Large\textbf{September 20XX\\}}
    % Module code 
    {\Large\textbf{your-moodle-code MSc Dissertation\\}}
    % Supervisor
    {\Large\textbf{Supervisor: XXX\\}}

    \vspace{3cm}

    {\textbf{Word count: XXX\\
    Reference style: Nature\\}}

    \vspace{3cm}

    {\textbf{
    This dissertation is submitted in partial fulfilment of the requirements for
    the Master's degree in [your programme title], UCL.\\}}

    \vfill
\end{center}
\end{titlepage}

% ----------------------  ABSTRACT -----------------------
\newpage
\chapter*{Abstract} % the Asterix (*) indicates that this chapter will be added to the table of contents but no number will be present beside it.
\thispagestyle{fancy}

The role of an abstract varies depending on the context in which it is used, but in general, an abstract serves as a concise summary or overview of a larger document, such as a research paper, thesis, article, or presentation. Its primary purposes are to provide readers with a quick understanding of the main points and findings of the document and to help them decide whether they want to read the full text. 

\vspace{1cm}
\textbf{Keywords:} Keywords, Keywords, Keywords

% -----------------  DECLARATION  -------------------
\chapter*{Declaration}
\thispagestyle{fancy}

I have read and understood the College and Departmental statements and guidelines concerning plagiarism. I declare that:
\begin{itemize}
    \item This submission is entirely my own original work.
    \item Wherever published, unpublished, printed, electronic or other information sources have been used as a contribution or component of this work, these are explicitly, clearly and individually acknowledged by appropriate use of quotation marks, citations, references and statements in the text. It is XXX words in length.
\end{itemize}

% ------------------  TABLE OF CONTENTS --------------------
\newpage 
\tableofcontents 
\thispagestyle{fancy}

% -------------------  LIST OF FIGURES --------------------
% \pagenumbering{arabic} 
% \setcounter{page}{1}
\newpage 
%{\let\oldnumberline\numberline       
% Uncomment to add the word 'Figure' to figure number in List of Figures
%\renewcommand{\numberline}{\figurename~\oldnumberline}  
\listoffigures%}
\thispagestyle{fancy}
% Add List of figures into contents without any numeration 
\addcontentsline{toc}{chapter}{List of Figures}

% -------------------  LIST OF TABLES ---------------------
\newpage
\listoftables 
\thispagestyle{fancy}
% Add List of tables into contents without any numeration 
\addcontentsline{toc}{chapter}{List of Tables} 

% -----------------  LIST OF ABBREVIATIONS -------------------
\chapter*{List of abbreviations}
\thispagestyle{fancy}
\addcontentsline{toc}{chapter}{List of abbreviations}

% Abbreviations related to optimizer
\textbf{RGB :} red, green and blue



% -----------------  ACKNOWLEDGEMENTS  -------------------
\chapter*{Acknowledgements}
\thispagestyle{fancy}
\addcontentsline{toc}{chapter}{Acknowledgement}
Acknowledgments, often found at the beginning or end of a document or publication, serve several important roles depending on the context in which they are used. The primary purpose of acknowledgments is to express gratitude and recognition to individuals, organizations, or institutions that have contributed in various ways to the creation or support of the document. 

% ------------------  main body  --------------------
\include{chapter_1_Introduction}

\include{chapter_2_Literature_Review}

% -------------------  BIBLIOGRAPHY ---------------------
\newpage
\setstretch{1.0}
% \thispagestyle{fancy}
\printbibliography[title = {Bibliography}]
\thispagestyle{fancy}
\addcontentsline{toc}{chapter}{Bibliography} % Adds References chapter to Table of Contents


\end{document}
%  -------------------------------------------------
%  --------- The document ends from here ----------- 
%  -------------------------------------------------